\chapter{Elective Splenectomy}
\index{Elective Splenectomy}

\section*{Definition and Overview}
An elective splenectomy is a planned surgical operation to remove the spleen. The spleen is an organ in the upper left abdomen that filters the blood, removes old blood cells, fights infection, and stores platelets and white blood cells. Removing it is recommended when the spleen causes problems that cannot be managed otherwise, such as severe blood disorders, an enlarged spleen causing symptoms, or certain types of cyst or tumour. Because the spleen plays a role in fighting infection, people who have it removed need vaccinations and lifelong attention to infection prevention.

\section*{Epidemiology}
Elective splenectomy is performed much less often than in the past, because medical treatments for many of the conditions it once treated have improved, and because surgery preserves the spleen where possible. It is most commonly performed today for conditions such as immune thrombocytopenia (low platelets), hereditary spherocytosis, and symptomatic splenomegaly. In many countries, the operation is usually performed laparoscopically, with a short hospital stay and good outcomes.

\section*{Aetiology and Pathophysiology}
The spleen performs several vital functions, and removal is considered when these cause harm.
\begin{itemize}
    \item \textbf{Immune thrombocytopenia:} the spleen destroys platelets coated with antibodies, causing bleeding risk
    \item \textbf{Hereditary spherocytosis:} abnormally shaped red cells are destroyed by the spleen, causing anaemia and gallstones
    \item \textbf{Splenomegaly:} an enlarged spleen can cause pain, early fullness, low blood counts, and a risk of rupture
    \item \textbf{Haematological conditions:} some leukaemias, lymphomas, and myeloproliferative disorders
    \item \textbf{Cysts and tumours:} splenic cysts, abscesses, and benign or malignant tumours
    \item \textbf{Role of the spleen:} it filters blood, removes old cells, makes antibodies, and stores platelets
    \item \textbf{After removal:} other organs, particularly the liver, take over many of the spleen's functions, but protection against certain bacteria is reduced
\end{itemize}

\section*{Clinical Features}
The symptoms that lead to elective splenectomy depend on the underlying condition.
\begin{itemize}
    \item Pain or fullness in the upper left abdomen
    \item Easy bruising and bleeding from low platelets
    \item Fatigue and paleness from anaemia
    \item Recurrent infections from reduced immune function
    \item A palpable, enlarged spleen
    \item Low blood counts on testing
    \item No symptoms at all in some people whose condition is detected incidentally
\end{itemize}

\section*{Differential Diagnosis}
The underlying conditions requiring splenectomy must be carefully diagnosed.
\begin{itemize}
    \item Immune thrombocytopenia versus other causes of low platelets
    \item Hereditary spherocytosis versus other causes of anaemia
    \item Splenomegaly from portal hypertension, infection, or blood disease
    \item Benign versus malignant causes of splenic masses
    \item Conditions that can be treated medically rather than surgically
\end{itemize}

\section*{When to Seek Medical Care}
People with the conditions listed above should be managed by a specialist team. Anyone with a damaged or enlarged spleen should seek urgent care for fever or injury, and should carry information about their condition.

\textbf{Contact your local emergency services for:}
\begin{itemize}
    \item Sudden severe left upper abdominal pain, which may indicate splenic rupture
    \item Fever and feeling unwell in someone without a functioning spleen
    \item Sudden severe bleeding or bruising
    \item Collapse or signs of sepsis
\end{itemize}

\section*{Diagnosis and Investigations}
Before surgery, the diagnosis and the fitness for surgery are assessed.
\begin{itemize}
    \item \textbf{Blood tests:} full blood count, blood film, and tests for the underlying condition
    \item \textbf{Imaging:} ultrasound, CT, or MRI of the spleen
    \item \textbf{Bone marrow biopsy:} for some haematological conditions
    \item \textbf{Pre-operative assessment:} general health, medications, and anaesthetic review
    \item \textbf{Vaccination:} immunisation against encapsulated bacteria is given before surgery
    \item \textbf{Infection screening:} checks for conditions such as malaria or HIV that affect the spleen
\end{itemize}

\section*{Management}
Management involves preparation, the operation, and lifelong aftercare.
\begin{itemize}
    \item \textbf{Vaccinations before surgery:} pneumococcal, meningococcal, Haemophilus influenzae type b, and influenza vaccines
    \item \textbf{Laparoscopic splenectomy:} the usual approach, with smaller incisions and faster recovery
    \item \textbf{Open splenectomy:} used when the spleen is very large or surgery is complex
    \item \textbf{Partial splenectomy:} preserving part of the spleen in selected children and adults
    \item \textbf{Aftercare:} pain relief, early mobilisation, and wound care
    \item \textbf{Lifelong precautions:} prompt treatment of fever, daily antibiotics in some people, and booster vaccinations
\end{itemize}

\section*{Complications}
\begin{itemize}
    \item The general risks of surgery: bleeding, infection, and anaesthetic complications
    \item Injury to nearby organs, including the pancreas, stomach, and bowel
    \item Blood clots, particularly in the veins of the portal system
    \item Overwhelming infection from encapsulated bacteria, the most serious long-term risk
    \item Raised platelet counts after surgery
    \item The effects of the underlying blood condition, which may persist
\end{itemize}

\section*{Prognosis and Outcomes}
The outlook after elective splenectomy is very good. Most people recover fully, and the procedure often cures or greatly improves the underlying condition, such as immune thrombocytopenia or hereditary spherocytosis. The main long-term consideration is the small but real lifelong risk of serious infection, which is largely prevented by vaccination, antibiotics, and prompt treatment of illness.

\section*{Prevention and Self-Care}
\begin{itemize}
    \item Complete all recommended vaccinations before and after surgery
    \item Carry a card or wear a bracelet identifying that you have no functioning spleen
    \item Seek medical care promptly for any fever or infection
    \item Take prophylactic antibiotics exactly as prescribed
    \item Avoid contact with people with serious infections where possible
    \item Attend scheduled follow-up and blood tests
    \item Consider travel advice from your doctor, including extra vaccines
\end{itemize}

\section*{Sources and Further Reading}
The following reputable sources provide further information for healthcare students and patients seeking reliable, current advice about splenectomy.
\begin{itemize}
    \item Australian Government Department of Health. \textit{Post-splenectomy vaccination advice.} https://www.health.gov.au/
    \item NHS. \textit{Splenectomy: surgery and aftercare.} https://www.nhs.uk/conditions/splenectomy/
    \item Mayo Clinic. \textit{Splenectomy: what to expect.} https://www.mayoclinic.org/tests-procedures/splenectomy/
    \item Healthdirect Australia. \textit{Spleen and splenectomy.} https://www.healthdirect.gov.au/spleen
    \item MSD Manual. \textit{Asplenia and splenic dysfunction.} https://www.msdmanuals.com/
    \item MedlinePlus. \textit{Spleen removal.} https://medlineplus.gov/ency/article/002944.htm
\end{itemize}
